% ============================================================
% РАЗДЕЛ 4. P2P/FUSE-архитектура кэширования ML-моделей
% ============================================================
\section{РАЗРАБОТКА P2P/FUSE-СИСТЕМЫ КЭШИРОВАНИЯ ML-МОДЕЛЕЙ}

Данный раздел описывает основную разработку работы --- специализированную систему распределённого кэширования ML-моделей, устраняющую ограничения NFS-этапа. Система построена на трёх ключевых компонентах: Redis-реестре метаданных (Control Plane), агенте \texttt{storage-agent} на каждой worker-ноде (Data Plane) и FUSE-клиенте \texttt{storage-mounter}, предоставляющем ML-runtime стандартный файловый интерфейс. Вместе они реализуют трёхуровневую иерархию доступа к данным: локальный диск ноды, P2P-сеть кластера и внешний репозиторий как резервный источник.

\subsection{Общая идея и требования к архитектуре}

Центральная идея архитектуры состоит в том, что локальные диски worker-нод образуют распределённый кэш, доступный всем подам в кластере. Данные модели хранятся не в одном месте, а распределены между нодами в виде чанков --- блоков фиксированного размера. При запуске нового пода ML-runtime обращается к ближайшей копии данных, а не к централизованному NFS или внешнему репозиторию. Этот подход переносит хранилище данных из единого централизованного узла в локальность вычислений, что соответствует принципу data locality, давно утвердившемуся в распределённых системах обработки данных.

Переход от NFS-этапа к P2P/FUSE-архитектуре мотивирован фундаментальным ограничением сетевых файловых систем: единый NFS-сервер становится узким местом при одновременном запуске нескольких подов с тяжёлыми ML-моделями. Если 10 подов одновременно читают модель весом 70~ГБ с одного NFS-сервера, суммарный поток данных составляет сотни гигабайт, что превышает пропускную способность большинства сетевых интерфейсов. P2P-архитектура устраняет это ограничение: каждый под обращается к локальному диску своей ноды или к нескольким соседним нодам, распределяя нагрузку на чтение горизонтально.

Требования, которым должна удовлетворять архитектура, сформулированы исходя из ограничений предшествующих подходов:

\begin{itemize}[leftmargin=*]
  \item \textbf{Минимальная задержка первого байта}: ML-runtime должен получить первые данные модели в течение миллисекунд, не дожидаясь полной загрузки весов. Это обеспечивается механизмом lazy loading через FUSE.
  \item \textbf{Отсутствие единой точки отказа}: падение одной ноды не должно делать модель недоступной. Репликация чанков на несколько нод и fallback на внешний репозиторий обеспечивают это свойство.
  \item \textbf{Прозрачность для ML-runtime}: фреймворки vLLM и TGI не должны требовать модификации кода для работы с системой. FUSE предоставляет стандартный POSIX-интерфейс.
  \item \textbf{Масштабируемость}: добавление новых нод должно автоматически увеличивать совокупный кэш без перенастройки.
  \item \textbf{Устойчивость к stale metadata}: устаревшие записи об ушедших нодах должны автоматически удаляться через TTL механизм Redis.
  \item \textbf{Совместимость с Kubernetes}: компоненты должны разворачиваться стандартными Kubernetes-примитивами (DaemonSet, ConfigMap, Secret) без изменения прав хостовой системы, за исключением доступа к диску.
\end{itemize}

Разработанная архитектура логически делится на три уровня.

\textbf{Уровень 1. Control Plane (Реестр метаданных).} Роль центрального источника информации выполняет Redis-кластер. В нём хранятся метаданные о каждом чанке каждой модели: на каких нодах он расположен, каков его размер и хеш. Кроме того, каждый активный \texttt{storage-agent} периодически обновляет запись о себе (heartbeat) с TTL. Redis выступает «телефонной книгой» кластера: любой компонент может узнать, у каких агентов есть нужный чанк, за одну операцию чтения. Control Plane намеренно отделён от Data Plane: метаданные не смешиваются с данными, что упрощает масштабирование и отказоустойчивость каждого уровня независимо.

\textbf{Уровень 2. Data Plane --- локальный (Инференс и виртуальная файловая система).} На каждой compute-ноде работает контейнер с ML-движком и компонент \texttt{storage-mounter}. FUSE-клиент монтирует виртуальную файловую систему в локальный путь пода. Для ML-движка создаётся иллюзия, что файлы модели уже целиком лежат на диске. При чтении определённого диапазона байтов \texttt{storage-mounter} перехватывает системный вызов и по Unix Domain Socket запрашивает данные у локального \texttt{storage-agent}. UDS выбран намеренно: коммуникация остаётся в пределах одного хоста и не нагружает сеть датацентра.

\textbf{Уровень 3. Data Plane --- горизонтальный (P2P-сеть).} При промахе локального кэша \texttt{storage-agent} обращается к Redis для поиска узла, имеющего нужный чанк. Данные передаются напрямую по TCP между агентами на разных нодах. Если чанк отсутствует во всём кластере --- выполняется fallback-загрузка из внешнего репозитория (HuggingFace, Ollama, ModelScope). Горизонтальный Data Plane реализует свойство self-healing кэша: каждая загрузка из внешнего источника наполняет локальный кэш, уменьшая вероятность повторного обращения к внешнему репозиторию.

\begin{figure}[htbp]
\centering
\includegraphics[width=0.94\textwidth]{arch_target.png}
\caption{Целевая P2P/FUSE-архитектура распределённого кэширования ML-моделей}
\label{fig:arch-target-p2p-fuse}
\end{figure}

Рисунок~\ref{fig:arch-target-p2p-fuse} визуализирует полный маршрут чтения данных в предложенной системе. Локальные обращения \texttt{storage-mounter} к \texttt{storage-agent} изолированы в пределах ноды через UDS, межузловой обмен чанками выполняется напрямую между агентами, а Redis используется только как реестр метаданных и не включается в тракт передачи самих весов. Это принципиально отличает предложенный подход от NFS: центральный компонент сохраняется лишь для координации, тогда как основной поток чтения распределяется между рабочими узлами.

\subsection{Модель данных: репозиторий, модель, ревизия, файл, чанк}

Для корректного кэширования и адресации данных принята иерархическая модель, отражающая структуру реестров ML-моделей. Иерархия адресации включает следующие уровни:

\begin{itemize}[leftmargin=*]
  \item \textbf{registry} --- источник модели: \texttt{huggingface}, \texttt{ollama}, \texttt{modelscope};
  \item \textbf{repository} --- пространство имён владельца (например, \texttt{deepseek-ai});
  \item \textbf{model} --- имя модели (например, \texttt{DeepSeek-LLM-67B-Chat});
  \item \textbf{revision} --- версия модели (commit hash или тег); обеспечивает иммутабельность кэша;
  \item \textbf{filepath} --- относительный путь к файлу внутри репозитория (например, \texttt{model-00001-of-00004.safetensors});
  \item \textbf{chunk\_id} --- порядковый номер чанка в файле (начиная с 0);
  \item \textbf{chunk\_size} --- фиксированный размер чанка (настраивается, типично 8--64~МБ).
\end{itemize}

Иммутабельность через revision принципиально важна для корректности кэша. Файлы конкретной ревизии модели никогда не изменяются: если revision зафиксирован как конкретный commit hash, то содержимое файлов при данном hash всегда одинаково. Поэтому кэшированный чанк всегда корректен, пока revision в ключе совпадает с запрошенной версией. Это устраняет необходимость в механизмах инвалидации кэша при обновлении модели: новая версия получает новый revision, и её чанки хранятся независимо от чанков предыдущей версии. Старые чанки вытесняются по LRU-политике по мере заполнения диска.

Физический путь к чанку на диске агента строится следующим образом:

\begin{lstlisting}[numbers=none,frame=none,basicstyle=\small\ttfamily]
<cache_dir>/
  <registry>/<repo>/<model>/<revision>/
    <hash(filepath)[0:2]>/
      <hash(filepath)[2:4]>/
        <hash(filepath)>/
          chunk_<chunk_id>
\end{lstlisting}

Хеширование пути к файлу решает две задачи. Первая --- устранение проблемы специальных символов и длинных имён: имена файлов в HuggingFace-репозиториях могут содержать слеши, пробелы и длинные пути, несовместимые с прямым использованием в качестве имени файла на диске. Хеш SHA-256 от filepath даёт детерминированное, фиксированной длины имя директории. Вторая задача --- равномерное распределение файлов: файловые системы Linux имеют ограничения на число записей в одной директории (\texttt{ext4} --- до нескольких десятков тысяч записей без \texttt{dir\_index}, что вызывает деградацию производительности); двухуровневое шардирование по первым четырём символам хеша обеспечивает $256 \times 256 = 65536$ возможных поддиректорий, практически исключая перегрузку отдельных директорий.

Ключ в Redis для записи о чанке строится аналогично иерархии на диске:

\begin{lstlisting}[numbers=none,frame=none,basicstyle=\small\ttfamily]
<registry>:<repo>:<model>:<revision>:<hash(filepath)>:<chunk_id>
\end{lstlisting}

Такая схема позволяет находить все чанки конкретного файла по префиксу ключа через \texttt{SCAN}, а также атомарно получать список нод, хранящих конкретный чанк, одной командой \texttt{SMEMBERS}. Единообразие схемы адресации между диском и Redis упрощает отладку: по ключу в Redis можно немедленно определить физический путь на диске агента.

Размер чанка влияет на компромисс между накладными расходами и гранулярностью. Мелкие чанки (4--8~МБ) допускают более точное управление тем, какие части модели кэшируются, но увеличивают число метаданных в Redis и число операций FUSE. Крупные чанки (64--128~МБ) снижают накладные расходы, но могут приводить к загрузке данных, которые ML-runtime не запросит. Диапазон 16--32~МБ является практически обоснованным для последовательного чтения больших \texttt{.safetensors}-файлов в ML-инференсе.

\subsection{Redis как реестр метаданных чанков}

\subsubsection*{Обоснование выбора Redis}

Выбор Redis в роли Control Plane обусловлен рядом характеристик, критичных для задачи реестра метаданных в динамической P2P-среде.

\textbf{Суб-миллисекундная задержка.} Redis хранит данные в оперативной памяти и обеспечивает время отклика на операции чтения и записи, измеряемое микросекундами \cite{redis_cluster_docs}. При обработке запроса чанка \texttt{storage-agent} выполняет один \texttt{SMEMBERS}-запрос к Redis перед обращением к соседнему узлу: накладные расходы на этот lookup должны быть пренебрежимо малы по сравнению со временем сетевой передачи чанка (десятки--сотни миллисекунд). Суб-миллисекундная задержка Redis делает его не узким местом системы.

\textbf{Атомарность операций.} Команды \texttt{SADD}, \texttt{HSET}, \texttt{EXPIRE} выполняются атомарно. Это позволяет агенту регистрировать чанк без состояний гонки: добавить свой идентификатор в множество пиров (\texttt{SADD chunk:...:peers <node\_id>}) и установить метаданные (\texttt{HSET chunk:...:meta size ... checksum ...}) атомарно. Другой агент, запрашивающий \texttt{SMEMBERS chunk:...:peers}, получит либо запись целиком, либо вообще не получит идентификатор данного узла --- промежуточного состояния не существует.

\textbf{Встроенный TTL.} Каждый ключ в Redis может иметь ограниченное время жизни. Если агент прекратил публиковать heartbeat (нода упала или перезагрузилась), его записи исчезнут автоматически через TTL. Это устраняет необходимость в явном механизме garbage collection: реестр самоочищается без дополнительного координатора. Redis использует комбинацию ленивого удаления (при обращении к истёкшему ключу) и периодической активной очистки, обеспечивая эффективное управление памятью при наличии TTL.

\textbf{Горизонтальное масштабирование.} Redis Cluster шардирует ключи по 16384 хеш-слотам, распределяя нагрузку между мастер-узлами. При отказе мастера кластер автоматически повышает реплику \cite{redis_cluster_docs}. Redis Sentinel решает аналогичную задачу для меньших кластеров: мониторит мастер, выполняет автоматическое переключение на реплику при его недоступности \cite{redis_sentinel_docs}. Оба решения устраняют Redis как единую точку отказа, что было ключевым недостатком NFS-этапа.

\textbf{Сравнение с альтернативами.} etcd обеспечивает строгую линейную консистентность через алгоритм Raft, но вносит задержку в единицы миллисекунд и оптимизирован для небольшого объёма редко меняющихся данных конфигурации --- не для высокочастотного реестра тысяч чанков. ZooKeeper не имеет встроенного TTL и требует сложного клиентского кода для имитации его функциональности; кроме того, ZooKeeper написан на Java и сложен в операционном обслуживании. Consul избыточен: большинство его возможностей (service mesh, ACL, health checks) не используются при решении задачи реестра чанков. Redis обеспечивает точно необходимый набор: быстрое чтение/запись эфемерных метаданных, автоматический TTL, богатые структуры данных (Sets, Hashes), горизонтальное масштабирование.

\subsubsection*{Схема ключей}

В реестре применяется следующая трёхуровневая схема ключей:

\begin{itemize}[leftmargin=*]
  \item \texttt{node:<node\_id>} --- heartbeat агента. Значение (JSON или строка) содержит IP-адрес, TCP-порт, load score, метку времени последнего обновления. TTL = 30--60~с. Если агент не обновляет heartbeat, запись автоматически исчезает, и следующий запрос к этому агенту не выполняется.
  \item \texttt{chunk:<logical\_id>:peers} --- Set идентификаторов нод, хранящих данный чанк. Операция \texttt{SADD <node\_id>} регистрирует ноду; \texttt{SMEMBERS} возвращает всех известных пиров за $O(1)$ (одна хеш-таблица в памяти Redis). Сложность \texttt{SMEMBERS} для Set из $k$ пиров --- $O(k)$.
  \item Hash с метаданными чанка строится по следующей схеме:
\begin{lstlisting}[numbers=none,frame=none,basicstyle=\footnotesize\ttfamily]
model:<registry>:<repo>:<model>:<revision>:
  file:<hash>:chunk:<id>:meta
\end{lstlisting}
  В Hash хранятся размер в байтах (\texttt{size}) и контрольная сумма SHA-256 (\texttt{checksum}); ключ используется для верификации целостности после загрузки.
\end{itemize}

Разделение информации о пирах и метаданных чанка на два ключа (\texttt{:peers} и \texttt{:meta}) позволяет обновлять их независимо. Метаданные чанка (размер, хеш) не изменяются после первой записи --- ключ \texttt{:meta} фактически write-once. Множество пиров (\texttt{:peers}) динамически растёт по мере того, как новые ноды кэшируют чанк, и может уменьшаться при явном \texttt{SREM} или истечении TTL.

После успешного сохранения чанка на диске агент выполняет следующую атомарную последовательность: \texttt{SADD chunk:\ldots:peers <node\_id>} и \texttt{HSET chunk:\ldots:meta size <n> checksum <sha256>}. Это обеспечивает согласованность реестра с фактическим состоянием дисков.

\subsubsection*{Heartbeat и управление жизненным циклом}

Каждый \texttt{storage-agent} публикует heartbeat каждые 10--30~с командой:

\begin{verbatim}
SET node:<node_id> <metadata_json> EX <ttl>
\end{verbatim}

Команда \texttt{SET ... EX} атомарно записывает значение и устанавливает TTL. Это исключает ситуацию, когда значение записано, но TTL ещё не установлен --- состояние, которое возникло бы при раздельном вызове \texttt{SET} и \texttt{EXPIRE}.

TTL выбирается в диапазоне 2--3 интервалов heartbeat: при TTL = 60~с и интервале heartbeat = 20~с нода выдерживает 3 пропущенных heartbeat прежде чем её запись устареет. Это обеспечивает устойчивость к кратковременным сбоям сети (задержке heartbeat) без риска оставить в реестре записи реально отказавших нод.

Любой агент, запрашивающий список пиров для чанка, проверяет свежесть каждого кандидата: \texttt{EXISTS node:<peer\_id>}. Если heartbeat истёк, агент пропускает данный пир и переходит к следующему. Это двухуровневая проверка: TTL heartbeat --- первый уровень защиты, явная проверка \texttt{EXISTS} --- второй.

\subsection{Обнаружение узлов: Gossip, mDNS и Redis heartbeat}

Для построения реестра активных агентов рассматривались три подхода.

\subsubsection*{Gossip Protocol}

Gossip (эпидемический протокол) --- децентрализованный механизм, при котором каждый узел периодически обменивается состоянием с $k$ случайно выбранными соседями \cite{gossip_protocol_overview}. Выбор соседей случаен и независим в каждом раунде, что обеспечивает экспоненциальное распространение информации: после $r$ раундов информация достигает примерно $(k{+}1)^r$ узлов (при отсутствии пересечений). Для кластера из $N$ нод информация распространяется за:

\begin{equation}
T_{\mathrm{gossip}} = O(\log N).
\label{eq:gossip}
\end{equation}

Для кластера из 1000 нод это порядка 10 раундов; для 25\,000 нод --- около 15 раундов. Каждый раунд занимает время, равное интервалу gossip (обычно 1--5~с), что означает задержку сходимости от 10 до 75~с для крупных кластеров. Именно этот временной разрыв делает Gossip неприемлемым для задачи реестра чанков.

\textbf{Преимущества Gossip}: полная децентрализация без центрального координатора, высокая отказоустойчивость --- протокол продолжает работать при отказе произвольного числа узлов, масштабируемость до сотен тысяч нод. Gossip применяется в Redis Cluster для обмена информацией о конфигурации слотов и состоянии узлов, а также в Apache Cassandra и HashiCorp Consul для поддержания актуальной топологии. Для этих применений eventual consistency приемлема, поскольку топология меняется редко.

\textbf{Недостатки для задачи реестра чанков}: задержка сходимости O(log N) раундов неприемлема при cold start, где каждая миллисекунда критична. Gossip создаёт постоянный фоновый трафик (каждый узел периодически инициирует обмен с несколькими соседями), что нагружает сеть. Отладка eventual-consistent системы существенно сложнее: нет гарантии, что в момент запроса реестр отражает актуальное состояние. Наконец, Gossip не даёт детерминированного ответа на вопрос «где сейчас находится чанк X» --- только вероятностный.

\subsubsection*{mDNS/UDP Broadcast}

Multicast DNS (RFC 6762) позволяет устройствам в локальной сети анонсировать себя через UDP-мультикаст на адрес 224.0.0.251, порт 5353 \cite{mdns_k8s_limitations}. Каждое устройство периодически отправляет mDNS Announcement, сообщая своё имя хоста и IP-адрес. Другие устройства могут найти сервис через mDNS Query.

В качестве механизма zero-configuration networking mDNS прост в реализации и не требует дополнительной инфраструктуры --- достаточно стандартной поддержки в ОС. Именно это делает mDNS популярным в домашних сетях (Apple Bonjour, DNS-SD).

Однако mDNS имеет серьёзные ограничения в среде Kubernetes \cite{mdns_k8s_limitations}:
\begin{itemize}[leftmargin=*]
  \item \textbf{Overlay-сети блокируют мультикаст}: большинство CNI-плагинов (Calico, Flannel, Cilium) используют VXLAN или IPIP, которые инкапсулируют L2-трафик в L3-пакеты и не распространяют мультикаст-фреймы между хостами.
  \item \textbf{Отсутствие IGMP/PIM в облаке}: облачные провайдеры (AWS, GCP, Azure) не поддерживают multicast routing на уровне underlay-сети.
  \item \textbf{Нарушение сетевой изоляции}: обходной путь через \texttt{hostNetwork: true} даёт поду доступ к мультикасту хостовой сети, но нарушает принципы сетевой изоляции Kubernetes и создаёт риски безопасности.
  \item \textbf{Ограниченная масштабируемость}: мультикаст масштабируется только в пределах одного broadcast-домена (VLAN), что несовместимо с multi-zone кластерами, охватывающими несколько зон доступности.
\end{itemize}

Таким образом, mDNS неприменим для обнаружения узлов в типичном production Kubernetes-кластере.

\subsubsection*{Redis heartbeat: выбранный подход}

Каждый \texttt{storage-agent} периодически обновляет запись о себе в Redis с TTL. Поиск активных агентов сводится к одному \texttt{SCAN} по префиксу \texttt{node:*} или \texttt{SMEMBERS} для конкретного чанка. Это обеспечивает время поиска $O(1)$, что принципиально отличает Redis heartbeat от Gossip ($O(\log N)$ раундов).

\begin{table}[htbp]
\centering
\caption{Сравнение алгоритмов обнаружения узлов}
\label{tab:discovery}
\begin{tabularx}{\textwidth}{|>{\raggedright\arraybackslash}p{42mm}|>{\centering\arraybackslash}p{22mm}|>{\centering\arraybackslash}p{29mm}|>{\centering\arraybackslash}p{29mm}|}
\hline
\textbf{Критерий} & \textbf{Gossip} & \textbf{\shortstack{mDNS/\\Broadcast}} & \textbf{Redis heartbeat} \\ \hline
Скорость lookup      & $O(\log N)$ раундов & Быстро (локальная сеть) & $O(1)$ \\
Консистентность      & eventual              & Нет                     & TTL-based \\
Поддержка в K8s      & Да                   & Плохая                  & Да \\
SPOF                 & Нет                  & Нет                     & Redis (без HA) \\
Сложность реализации & Средняя              & Низкая                  & Низкая \\
Масштабируемость     & Высокая              & Ограниченная            & Высокая (Cluster) \\
Гарантия актуальности & Нет (eventual)     & Только локальная сеть   & TTL-ограниченная \\
\hline
\end{tabularx}
\end{table}

Таблица~\ref{tab:discovery} обобщает ключевые различия трёх подходов. Из анализа следует, что для кластеров Kubernetes умеренного масштаба (до нескольких тысяч нод) Redis heartbeat обеспечивает оптимальный баланс: минимальная задержка поиска, детерминированный ответ, прозрачная интеграция с Kubernetes-инфраструктурой.

Единственный риск --- зависимость от доступности Redis --- устраняется развёртыванием Redis Sentinel (для кластеров малого и среднего масштаба) или Redis Cluster с несколькими мастерами и репликами (для больших кластеров). При временной недоступности Redis \texttt{storage-agent} переключается на локальный кэш и fallback, сохраняя работоспособность.

\subsection{Компонент storage-agent: DaemonSet-агент на каждой ноде}

\texttt{storage-agent} --- агент, развёртываемый как DaemonSet на каждой worker-ноде кластера. Kubernetes DaemonSet гарантирует, что на каждой подходящей ноде запущен ровно один экземпляр, который автоматически создаётся при добавлении новой ноды и удаляется при её выводе из эксплуатации. Агент монтирует локальный диск ноды (через \texttt{hostPath} volume) и использует его как кэш-хранилище чанков.

Использование DaemonSet вместо Deployment критически важно: именно DaemonSet обеспечивает ковенант «один агент на ноду», без которого невозможно гарантировать, что локальный кэш ноды управляется ровно одним процессом. Конкурентный доступ нескольких агентов к одному кэш-директории потребовал бы механизма блокировок, что усложнило бы архитектуру.

Агент является основным компонентом Data Plane и выполняет следующие функции:

\begin{itemize}[leftmargin=*]
  \item управление локальным диском и chunk-кэшем ноды (запись, чтение, вытеснение по LRU);
  \item обслуживание запросов от локального \texttt{storage-mounter} через Unix Domain Socket (в целевой архитектуре);
  \item обслуживание TCP-запросов от агентов на соседних нодах (P2P-поставщик чанков);
  \item периодическая публикация heartbeat и метаданных чанков в Redis;
  \item fallback-загрузка чанков из внешнего репозитория при их отсутствии в кластере.
\end{itemize}

\subsubsection*{HTTP API и текущая реализация}

В текущем прототипе реализован HTTP-API с двумя основными эндпоинтами:

\begin{itemize}[leftmargin=*]
  \item получение дерева файлов модели (имена файлов, размеры, хеши):
\begin{lstlisting}[numbers=none,frame=none,basicstyle=\footnotesize\ttfamily,aboveskip=0pt,belowskip=0pt]
GET /api/v1/models/{registry}/{repo}/{model}/
  {revision}/tree
\end{lstlisting}
  \item чтение диапазона байтов файла модели:
\begin{lstlisting}[numbers=none,frame=none,basicstyle=\footnotesize\ttfamily,aboveskip=0pt,belowskip=0pt]
GET /api/v1/models/{registry}/{repo}/{model}/
  {revision}/raw/{filepath}
\end{lstlisting}
\end{itemize}

Диапазон байтов в запросе чтения задаётся HTTP-заголовком \texttt{Range: bytes=X-Y}.

Выбор HTTP как протокола в текущем прототипе обусловлен простотой реализации и отладки. В целевой архитектуре для локальной связи \texttt{storage-mounter $\to$ storage-agent} планируется переход на бинарный протокол поверх UDS для снижения накладных расходов.

Компонент \texttt{model-gateway} реализует приоритетный порядок источников при обработке запроса чанка:
\begin{enumerate}[leftmargin=*]
  \item проверяется наличие чанка в локальном хранилище (\texttt{localstorage}) через \texttt{os.Stat} по физическому пути;
  \item при промахе инициируется загрузка из HuggingFace через HTTP Range Request;
  \item загруженный чанк записывается на локальный диск атомарно: сначала во временный файл в той же файловой системе, затем переименовывается через \texttt{os.Rename}.
\end{enumerate}

Атомарная запись через \texttt{os.Rename} критически важна: системный вызов \texttt{rename(2)} на Linux атомарен при условии, что исходный и целевой пути находятся в одной файловой системе. Это гарантирует, что конкурентный читатель никогда не увидит частично записанный чанк: файл либо полностью существует, либо отсутствует. Этот инвариант устраняет необходимость в явных блокировках при параллельных запросах к одному чанку.

Компонент \texttt{localstorage} организует файлы на диске по иерархии, описанной в разделе~4.2. Доступ к метаданным файловой системы (размер чанка, наличие файла) осуществляется через \texttt{os.Stat}, что исключает необходимость отдельной локальной базы данных: файловая система сама является индексом локального кэша.

\subsubsection*{Поддержка singleflight для предотвращения thundering herd}

При одновременном поступлении нескольких запросов одного и того же чанка (например, когда несколько горутин в \texttt{storage-mounter} запрашивают один чанк параллельно) без специальных мер каждый запрос инициирует отдельную загрузку из внешнего источника. Это явление называется thundering herd. Для его предотвращения планируется использование механизма \texttt{golang.org/x/sync/singleflight}: если загрузка чанка уже выполняется, повторные запросы ожидают её завершения и получают тот же результат. Это сокращает нагрузку на внешний репозиторий и ускоряет ответ повторным запросам.

\subsubsection*{Целевое P2P-расширение}

В целевой архитектуре \texttt{storage-agent} дополняется логикой P2P-обмена. При промахе локального кэша агент выполняет следующую последовательность:
\begin{enumerate}[leftmargin=*]
  \item запрашивает Redis: \texttt{SMEMBERS chunk:<id>:peers} для получения списка нод с нужным чанком;
  \item фильтрует список: проверяет \texttt{EXISTS node:<peer\_id>} для каждого кандидата;
  \item ранжирует кандидатов по load score и задержке;
  \item скачивает чанк по TCP от выбранного пира;
  \item сохраняет на локальный диск атомарно через \texttt{os.Rename};
  \item регистрирует себя в Redis: \texttt{SADD chunk:<id>:peers <local\_node\_id>}.
\end{enumerate}

Этот процесс замыкает P2P-цикл: каждый новый потребитель чанка становится потенциальным поставщиком для следующего. После нескольких таких циклов чанки популярных моделей распространяются по кластеру органически, без центрального координатора.

\subsubsection*{Многоуровневый кэш storage-agent}

Эффективность \texttt{storage-agent} как прокси-шлюза между \texttt{storage-mounter} и внешними реестрами определяется организацией кэша. В прототипе реализована трёхуровневая иерархия, каждый уровень которой обслуживает свой класс запросов.

\textbf{L0~--- горячий уровень в оперативной памяти.} Компонент \texttt{inMemoryChunkCache} реализует SLRU-кэш (Segmented Least Recently Used) с ёмкостью, задаваемой в байтах. Внутри кэша поддерживается двусвязный список активных чанков и два хеш-словаря: первый сопоставляет ключ (идентификатор модели и номер чанка) с узлом списка, второй — с самими данными. При обращении к чанку его узел перемещается в начало списка. При переполнении ёмкости вытеснение выполняется с хвоста списка, пока не освободится достаточно места. Время операций Get и Set составляет $O(1)$. Уровень L0 обслуживает повторные обращения к одним и тем же чанкам без дискового ввода-вывода; при перезапуске агента его содержимое теряется, однако это допустимо, поскольку данные восстановимы с диска или из внешнего источника.

\textbf{L1~--- персистентный кэш на локальном диске.} Каждый чанк хранится как отдельный файл в иерархии директорий, описанной в разделе~4.2. Запись на диск выполняется по схеме write-behind: метод \texttt{Write()} помещает чанк в буферизованный канал и немедленно возвращает управление, не дожидаясь физической записи. Отдельная горутина-воркер читает из канала и записывает файл атомарно~--- сначала во временный файл в той же файловой системе, затем вызывает \texttt{os.Rename()}. Атомарность \texttt{rename(2)} гарантирует, что конкурентный читатель никогда не увидит частично записанного чанка. Разделение пути записи и пути чтения устраняет задержку дисковой операции из критического пути ответа клиенту.

\textbf{Выбор файлового кэша.} Альтернативы на основе встроенных хранилищ~--- BadgerDB, BoltDB, LevelDB~--- не подходят для чанков размером 16--128~МБ. LSM-деревья (BadgerDB, Pebble) оптимизированы для записи небольших значений: крупные значения уходят в Value Log, что порождает write amplification и фрагментацию. BoltDB хранит все данные в едином B-tree файле с эксклюзивной блокировкой при записи, что создаёт конкуренцию под нагрузкой. Файловый кэш лишён этих ограничений: операционная система обеспечивает параллельный доступ к независимым файлам, а её page cache автоматически удерживает горячие чанки в RAM без дополнительного уровня буферизации.

\textbf{Расширение: SQLite-индекс для структурированного вытеснения.} В текущей реализации удаление устаревших чанков требует обхода файловой системы по времени последнего доступа (\texttt{atime}). Запланированное расширение --- таблица SQLite в WAL-режиме с записями \texttt{(model\_ref, chunk\_id, size\_bytes, last\_access\_unix)}~--- позволит выполнять LRU-вытеснение запросом вида \texttt{ORDER BY last\_access\_unix ASC LIMIT N}, не читая файловую систему. Пара SQLite WAL + файловый кэш сохраняет преимущества прямого I/O и при этом добавляет структурированный учёт использования.

\textbf{Улучшение: ARC вместо LRU для горячего уровня.} При смешанной нагрузке~--- повторные обращения к популярным моделям параллельно с разовыми сканированиями при первой загрузке~--- чистый LRU вытесняет горячие чанки разовыми промахами. Алгоритм ARC (Adaptive Replacement Cache) поддерживает два рабочих списка (T1~--- недавно добавленные, T2~--- повторно запрошенные) и два ghost-списка ключей (B1 и B2). Попадания в B1 сдвигают внутреннюю границу $p$ в пользу рецентности, в B2~--- в пользу частотности. Это делает ARC самонастраивающимся без ручного подбора параметров.

\subsection{Компонент storage-mounter: FUSE-клиент}

\texttt{storage-mounter} --- FUSE-клиент, предоставляющий ML-runtime стандартный файловый интерфейс для доступа к весам модели. Он монтирует виртуальную файловую систему в путь внутри пода (например, \texttt{/mnt/models}).

\subsubsection*{Механизм FUSE}

FUSE (Filesystem in Userspace) позволяет реализовать файловую систему в виде обычного пользовательского процесса без привилегий root и без изменения кода ядра \cite{fuse_kernel_docs}. Архитектура включает три взаимодействующих компонента:

\begin{itemize}[leftmargin=*]
  \item \textbf{Ядерный модуль \texttt{fuse.ko}}: перехватывает обращения VFS к FUSE-точке монтирования и передаёт запросы пользовательскому демону через символьное устройство \texttt{/dev/fuse}.
  \item \textbf{VFS (Virtual Filesystem Switch)}: единый интерфейс ядра Linux для всех файловых систем; именно он определяет, что файл принадлежит FUSE-ФС, и направляет запрос соответствующему обработчику.
  \item \textbf{Пользовательский демон}: реализует логику файловой системы; в данном случае --- \texttt{storage-mounter}, написанный на Go с использованием библиотеки \texttt{github.com/hanwen/go-fuse/v2} \cite{gofuse}.
\end{itemize}

Жизненный цикл системного вызова \texttt{read()} в FUSE выглядит следующим образом:
\begin{enumerate}[leftmargin=*]
  \item Приложение ML-runtime (vLLM) вызывает \texttt{read(fd, buf, count)}.
  \item Ядро перехватывает вызов через VFS, определяет принадлежность файла к FUSE-ФС.
  \item Ядерный модуль \texttt{fuse.ko} упаковывает запрос в пакет протокола FUSE и помещает в очередь.
  \item Пользовательский демон \texttt{storage-mounter} читает пакет из \texttt{/dev/fuse}.
  \item Демон обрабатывает запрос: вычисляет chunk\_id, обращается к \texttt{storage-agent}.
  \item \texttt{storage-agent} возвращает данные чанка.
  \item Демон записывает ответ обратно в \texttt{/dev/fuse}.
  \item Ядро передаёт данные приложению.
\end{enumerate}

Для ML-runtime этот процесс полностью прозрачен: он воспринимает FUSE-ФС как обычную локальную файловую систему. Это ключевое свойство позволяет интегрировать систему без изменения кода vLLM или TGI.

Преимущество FUSE перед альтернативами --- CSI volume с полным копированием весов в emptyDir --- в поддержке lazy loading: данные загружаются только при обращении к конкретному диапазону байтов. ML-runtime может начать инициализацию первых слоёв модели ещё до того, как остальные веса получены с диска или из сети. Это перекрывает фазы вычисления и ввода-вывода, сокращая время cold start.

\textbf{Накладные расходы FUSE.} Каждый системный вызов, направленный FUSE-ФС, требует переключения контекста kernel $\leftrightarrow$ userspace. Для рабочей нагрузки ML-инференса, при которой читаются большие последовательные блоки (\texttt{.safetensors}-файлы по несколько гигабайт), накладные расходы минимизируются за счёт крупных чанков: каждое переключение контекста амортизируется на несколько мегабайт полезных данных. Для мелкогранулярного случайного доступа FUSE может быть неоптимален, однако паттерн чтения ML-весов --- преимущественно последовательный.

\subsubsection*{Ключевые операции}

\texttt{storage-mounter} реализует следующие FUSE-операции \cite{gofuse}:

\begin{itemize}[leftmargin=*]
  \item \texttt{getattr} --- возвращает метаданные файла: размер, права, тип. Позволяет ML-runtime видеть полный размер файла модели без его загрузки. Вызывается при каждом обращении к пути и при выполнении \texttt{stat(2)}.
  \item \texttt{readdir} --- возвращает список файлов директории модели. Получается из HuggingFace API (\texttt{/api/models/<repo>/tree/<revision>}) и кэшируется локально на стороне клиента.
  \item \texttt{lookup} --- разрешает имя файла в inode. Вызывается при открытии файла по имени; возвращает атрибуты файла или ENOENT.
  \item \texttt{open} --- открывает файл для чтения. Может проверять права доступа и инициализировать внутренние структуры для отслеживания паттерна чтения.
  \item \texttt{read(offset, size)} --- читает диапазон байтов файла. Здесь происходит ключевая трансляция в chunk-запросы к \texttt{storage-agent}.
\end{itemize}

\subsubsection*{Трансляция read(offset, size) в chunk-запросы}

Ключевая операция \texttt{read(offset, size)} транслируется в запросы чанков следующим образом. По значениям \texttt{offset} и \texttt{chunk\_size} вычисляется диапазон затронутых чанков:

\begin{equation}
chunk\_id_{\mathrm{start}} = \left\lfloor \frac{offset}{chunk\_size} \right\rfloor, \quad
chunk\_id_{\mathrm{end}} = \left\lfloor \frac{offset + size - 1}{chunk\_size} \right\rfloor.
\label{eq:chunk_range}
\end{equation}

Для каждого чанка в диапазоне $[chunk\_id_{\mathrm{start}},\, chunk\_id_{\mathrm{end}}]$ отправляется отдельный запрос к \texttt{storage-agent}. При $chunk\_size = 16$~МБ и запросе на чтение 64~МБ, начиная с произвольного смещения, затронуто не более 5 чанков (4 полных + 1 граничный). Полученные данные объединяются в буфере, из которого вырезается точный запрошенный диапазон. Граничные чанки используются частично: первый --- начиная с \texttt{offset \% chunk\_size}, последний --- до \texttt{(offset + size - 1) \% chunk\_size + 1}.

\subsubsection*{Текущая реализация и целевое развитие}

В текущем прототипе \texttt{storage-mounter} реализован с использованием библиотеки \texttt{github.com/hanwen/go-fuse/v2} \cite{gofuse}. Поддержаны операции \texttt{Lookup}, \texttt{Getattr}, \texttt{Readdir}. Список файлов модели получается из HuggingFace API. Операция \texttt{read} в текущей реализации делегируется непосредственно к API HuggingFace через HTTP Range Request.

В целевой архитектуре операция \texttt{read} транслируется в запросы к локальному \texttt{storage-agent} через UDS: вместо прямого обращения к внешнему репозиторию \texttt{storage-mounter} запрашивает чанк у агента, который обслуживает три сценария --- local cache hit, peer hit, remote fallback.

Дополнительно планируется реализация prefetching: при обнаружении последовательного паттерна чтения \texttt{storage-mounter} асинхронно запрашивает следующие $N$ чанков до их явного запроса ML-runtime. Это позволяет скрыть задержку сетевой или дисковой передачи за временем обработки текущего чанка GPU-ядром.

\subsubsection*{Кэш и буферизация чтения в storage-mounter}

\texttt{storage-mounter} обрабатывает поток системных вызовов \texttt{read()} от ML-runtime и транслирует их в запросы чанков к \texttt{storage-agent}. Для минимизации задержки и нагрузки на агента в стороне клиента реализовано несколько механизмов буферизации.

\textbf{Кольцевой буфер (streaming ring buffer).} Компонент \texttt{streaming.Cache} хранит массив фиксированной длины \texttt{capacity}, в котором чанк с номером $k$ размещается по индексу $k \bmod \mathtt{capacity}$. Запись нового чанка всегда успешна и выполняется за $O(1)$: если в слоте уже находится другой чанк, он вытесняется без дополнительных сравнений. Чтение проверяет соответствие \texttt{chunk.Id == chunkId}: при коллизии двух чанков с одинаковым индексом результат возвращается как промах. Такая схема оптимальна для последовательного чтения одного файла: при линейном обходе каждый слот занимается следующим чанком ровно тогда, когда предыдущий уже не нужен.

\textbf{Предзагрузчик (prefetcher).} При каждом вызове \texttt{Read(chunkId)} предзагрузчик запускает до \texttt{prefetchAhead-1} горутин, каждая из которых загружает следующий чанк, если его нет в кольцевом буфере. Количество одновременно выполняемых предзагрузок ограничивается параметром \texttt{prefetchParallelism}, что предотвращает насыщение полосы пропускания UDS-соединения. Для исключения дублирующихся загрузок одного чанка несколькими горутинами применяется механизм \texttt{singleflight.Group}: если загрузка данного \texttt{chunkId} уже выполняется, повторный запрос блокируется и получает тот же результат без отдельного обращения к агенту. Эффект предзагрузки~--- перекрытие фаз вычисления GPU и сетевого ввода-вывода: пока GPU обрабатывает текущий слой модели, следующие чанки уже загружаются в буфер.

\textbf{Пул байтовых буферов.} Каждый чанк требует выделения байтового буфера объёмом \texttt{chunkSize}. При параллельных запросах нескольких FUSE-горутин без механизма переиспользования каждый вызов \texttt{make([]byte, chunkSize)} создаёт новую аллокацию: при \texttt{chunkSize = 32}~МБ и 64~параллельных горутинах это даёт 2~ГБ аллокаций в секунду, что насыщает сборщик мусора. Компонент \texttt{pool.Pool} реализует мультиразмерный пул через \texttt{sync.Pool}: для каждого уникального значения \texttt{chunkSize} хранится отдельный пул. Метод \texttt{Get(size)} возвращает переиспользованный буфер, \texttt{Put(buf)} возвращает его в пул после обработки чанка.

\textbf{Ring buffer vs per-file LRU.} Рабочая нагрузка ML-инференса преимущественно последовательна: слои \texttt{.safetensors}-файла читаются от начала к концу, поэтому кольцевой буфер покрывает этот сценарий без коллизий. Однако шардированные модели (несколько \texttt{.safetensors}-файлов, загружаемых одновременно) порождают коллизии: чанки разных файлов с одинаковым номером вытесняют друг друга по модулю ёмкости буфера. В этом сценарии предпочтительна per-file LRU-структура: отдельный LRU-кэш для каждого открытого файла изолирует их чанки. Реализация per-file LRU хранит словарь \texttt{filepath → lruCache} с очисткой записи при закрытии файла. Кольцевой буфер остаётся предпочтительным для одиночных больших файлов как более простой и эффективный; переход к per-file LRU оправдан при наличии двух и более одновременно открытых файлов с перекрывающимися диапазонами \texttt{chunkId}.

\subsection{UDS и TCP: разделение локальной и межнодовой коммуникации}

Архитектура применяет два различных транспортных протокола в зависимости от физической топологии взаимодействия. Разделение транспортных протоколов по границе «внутри ноды / между нодами» отражает фундаментальные различия в задержке и модели безопасности.

\subsubsection*{Unix Domain Socket: внутриузловая связь}

Unix Domain Socket (UDS) используется для связи \texttt{storage-mounter $\to$ local storage-agent} на одной ноде. UDS работает через файловую систему (\texttt{/var/run/storage-agent.sock}), минуя IP-стек полностью. Ключевые технические свойства UDS:

\begin{itemize}[leftmargin=*]
  \item \textbf{Отсутствие сетевого стека}: данные не инкапсулируются в IP/TCP-пакеты; нет расчёта контрольных сумм, маршрутизации, обработки заголовков, трёхстороннего рукопожатия.
  \item \textbf{Копирование через память ядра}: передача данных осуществляется прямым копированием между адресными пространствами процессов в ядре. Типичная задержка --- 1--5~мкс, пропускная способность --- до 10~ГБ/с.
  \item \textbf{Изоляция через права доступа}: права на socket-файл (\texttt{chmod/chown}) определяют, какие процессы могут устанавливать соединение. Механизм \texttt{SO\_PEERCRED} позволяет серверу получить PID, UID и GID клиента непосредственно от ядра без дополнительной аутентификации.
  \item \textbf{Нет NetworkPolicy}: UDS не пересекает сетевые границы Kubernetes, поэтому политики \texttt{NetworkPolicy} к нему неприменимы.
\end{itemize}

\subsubsection*{TCP: межнодовая связь}

TCP используется для связи \texttt{storage-agent $\to$ storage-agent} между нодами. TCP обеспечивает надёжную упорядоченную доставку данных, управление потоком и контроль перегрузки --- свойства, необходимые при передаче чанков объёмом 8--64~МБ по реальной сети. TCP-соединения совместимы со всеми CNI-плагинами и политиками \texttt{NetworkPolicy}.

Контроль перегрузки TCP (алгоритмы CUBIC, BBR) автоматически адаптирует скорость передачи к доступной пропускной способности канала и реагирует на потери пакетов. Это особенно важно при одновременной P2P-передаче нескольких чанков между несколькими парами нод: TCP предотвращает перегрузку коммутаторов датацентра.

\begin{table}[htbp]
\centering
\caption{Сравнение UDS и TCP в контексте системы}
\label{tab:uds_tcp}
\begin{tabularx}{\textwidth}{|>{\raggedright\arraybackslash}p{45mm}|>{\centering\arraybackslash}p{35mm}|>{\centering\arraybackslash}X|}
\hline
\textbf{Критерий} & \textbf{UDS} & \textbf{TCP} \\ \hline
Область применения     & Один узел                & Разные узлы \\
IP-стек                & Отсутствует              & Используется \\
Задержка               & 1--5~мкс                 & 0,1--1~мс (RTT в ДЦ) \\
Накладные расходы      & Минимальные              & Заголовки, handshake, CRC \\
Надёжность             & Да (ядро ОС)             & Да (ACK, retransmit) \\
Flow control           & Нет                      & Да (TCP window) \\
Congestion control     & Нет                      & Да (CUBIC, BBR) \\
Безопасность           & FS perms, SO\_PEERCRED   & TLS/mTLS (целевое) \\
K8s NetworkPolicy      & Не применима             & Применима \\
\hline
\end{tabularx}
\end{table}

Таблица~\ref{tab:uds_tcp} демонстрирует, что выбор протокола строго определяется физическим расположением компонентов: UDS --- когда оба процесса на одном хосте, TCP --- при межнодовом взаимодействии. Комбинация двух протоколов даёт оптимум: внутри ноды минимальная задержка UDS, между нодами надёжность и совместимость TCP.

\subsubsection*{Оценка пропускной способности через BDP}

Произведение полосы пропускания и задержки (BDP, Bandwidth-Delay Product) для TCP-канала:

\begin{equation}
\mathrm{BDP} = B \cdot \mathrm{RTT},
\label{eq:bdp}
\end{equation}

определяет объём данных, одновременно находящихся «в полёте» при полной утилизации канала. Для внутренней сети датацентра с $B = 10$~Гбит/с и $\mathrm{RTT} = 0.5$~мс:
\[
\mathrm{BDP} = \frac{10 \cdot 10^9}{8}\,\text{Б/с} \times 0.5 \cdot 10^{-3}\,\text{с} \approx 625\,\text{КБ}.
\]

Это означает, что для полной утилизации 10G-канала размер TCP-окна должен быть не менее 625~КБ. Типичный размер чанка 8--64~МБ существенно превышает BDP, что позволяет одному TCP-соединению утилизировать канал без специальной настройки размера окна \cite{hennessy2011ca}. При переходе на 25G и 100G сети BDP возрастает пропорционально, и может потребоваться явная настройка \texttt{net.core.rmem\_max} и \texttt{net.ipv4.tcp\_rmem}.

\subsection{Алгоритм чтения чанка}

\subsubsection*{Полная тринадцатишаговая последовательность}

Полный сценарий чтения чанка включает следующую последовательность шагов, охватывающую все три уровня архитектуры:

\begin{enumerate}[leftmargin=*]
  \item ML-runtime (vLLM) вызывает \texttt{read(filepath, offset, size)}.
  \item Ядро Linux перехватывает вызов через VFS и передаёт \texttt{storage-mounter} через \texttt{/dev/fuse}.
  \item \texttt{storage-mounter} вычисляет диапазон чанков по формуле~(\ref{eq:chunk_range}).
  \item \texttt{storage-mounter} отправляет запрос \texttt{(chunk\_id)} локальному \texttt{storage-agent} через UDS.
  \item \texttt{storage-agent} проверяет локальное хранилище операцией \texttt{os.Stat} по физическому пути.
  \item \textbf{Сценарий A (local cache hit)}: чанк найден на диске --- данные читаются из файла и немедленно возвращаются через UDS.
  \item \textbf{Сценарий B (local miss)}: агент запрашивает Redis: \texttt{SMEMBERS chunk:<id>:peers}.
  \item Агент проверяет свежесть каждого кандидата: \texttt{EXISTS node:<peer\_id>}.
  \item \textbf{Сценарий B.1 (peer cache hit)}: выбирается пир с минимальным score; устанавливается TCP-соединение; чанк загружается; сохраняется атомарно на диск через \texttt{os.Rename}; агент регистрирует себя: \texttt{SADD chunk:<id>:peers <local\_node>}.
  \item \textbf{Сценарий B.2 (peer miss, remote fallback)}: список пиров пуст или все пиры недоступны; агент выполняет HTTP Range Request к HuggingFace; полученный чанк сохраняется и регистрируется в Redis.
  \item Чанк возвращается \texttt{storage-mounter} через UDS.
  \item \texttt{storage-mounter} собирает данные всех запрошенных чанков и вырезает нужный диапазон.
  \item Данные возвращаются через FUSE в ML-runtime.
\end{enumerate}

\subsubsection*{Модель среднего времени доступа (AMAT)}

Среднее время доступа к чанку описывается трёхуровневой моделью AMAT (Average Memory Access Time) \cite{hennessy2011ca}:

\begin{equation}
\mathrm{AMAT} = T_{\mathrm{local}} + (1-H_l)\left[T_{\mathrm{p2p}} + (1-H_p) \cdot T_{\mathrm{remote}}\right],
\label{eq:amat}
\end{equation}

где:
\begin{itemize}[leftmargin=*]
  \item $T_{\mathrm{local}}$ --- время чтения чанка из локального диска (сценарий A);
  \item $H_l$ --- вероятность попадания в локальный кэш (local hit rate);
  \item $T_{\mathrm{p2p}}$ --- время получения чанка от соседнего агента по TCP (сценарий B.1);
  \item $H_p$ --- вероятность попадания в кэш соседних нод (peer hit rate);
  \item $T_{\mathrm{remote}}$ --- время загрузки чанка из внешнего репозитория (сценарий B.2).
\end{itemize}

Вывод формулы~(\ref{eq:amat}) основан на формуле полного математического ожидания. С вероятностью $H_l$ запрос обслуживается локально за время $T_{\mathrm{local}}$. С вероятностью $(1-H_l)$ происходит промах: система переходит на второй уровень, где с вероятностью $H_p$ чанк находится у соседнего агента ($T_{\mathrm{p2p}}$), и с вероятностью $(1-H_p)$ требуется загрузка из внешнего источника ($T_{\mathrm{remote}}$) \cite{li2020caching}. Перемножение условных вероятностей даёт коэффициенты перед каждым слагаемым.

После прогрева кластера (первая загрузка популярной модели завершена) можно ожидать $H_l \approx 0.85$ и $H_p \approx 0.90$. Тогда доля запросов, достигающих внешнего источника:
\[
(1-H_l)(1-H_p) = 0.15 \times 0.10 = 0.015.
\]
Около 1.5\% запросов уходят во внешний репозиторий, тогда как 98.5\% обслуживаются из локального или соседнего кэша. Это кардинально отличается от схемы \texttt{storage-initializer}, при которой каждый новый под загружает модель целиком из внешнего источника. Поскольку $T_{\mathrm{remote}}$ может составлять десятки секунд для модели весом 70~ГБ, снижение его коэффициента с 1.0 до 0.015 является определяющим фактором ускорения cold start.

\subsection{Политика кэширования и вытеснения}

\subsubsection*{Распределение популярности моделей: закон Ципфа}

Популярность моделей в реальных производственных системах распределена неравномерно: небольшое число наиболее востребованных моделей обеспечивает большую часть запросов. Аналогичная закономерность наблюдается в CDN (Content Delivery Networks), веб-кэшировании и системах хранения медиаконтента. Это распределение соответствует закону Ципфа \cite{zipf_law}:

\begin{equation}
P(k) = \frac{k^{-\alpha}}{\sum_{i=1}^{M} i^{-\alpha}},
\label{eq:zipf}
\end{equation}

где $k$ --- ранг модели по популярности (чем ниже ранг, тем популярнее), $\alpha$ --- параметр наклона (для ML-моделей оценивается в диапазоне $\alpha \approx 0.8$--$1.2$ по публичным оценкам исследований веб-трафика), $M$ --- общее число моделей.

Закон Ципфа означает, что модель с рангом 1 запрашивается примерно в $2^\alpha$ раз чаще модели с рангом 2. При $\alpha = 1$ и $M = 100$ моделях: топ-10 моделей (10\%) генерируют более 60\% запросов, топ-20 моделей --- около 75\%. Следствие для кэширования: даже кэш, вмещающий 20\% самых популярных моделей, обеспечивает hit rate выше 75\%. Нет необходимости хранить все существующие модели --- достаточно приоритизировать наиболее запрашиваемые.

\subsubsection*{LRU как политика вытеснения}

При превышении лимита дискового пространства применяется алгоритм LRU (Least Recently Used): вытесняются чанки, к которым дольше всего не было обращений. LRU хорошо согласуется с законом Ципфа: часто используемые модели (низкий ранг) регулярно обращаются к своим чанкам, обновляя метку «последнего использования», тогда как редко используемые (высокий ранг) постепенно вытесняются.

Практическая реализация LRU для файлового кэша: при обращении к чанку обновляется атрибут \texttt{atime} файла на диске (\texttt{O\_NOATIME} отключён намеренно для отслеживания времени доступа). При необходимости вытеснения перебираются файлы в директории кэша, сортируются по \texttt{atime} по возрастанию, и удаляются до восстановления свободного места.

\subsubsection*{Аппроксимация Чи для оценки hit rate}

Для приближённой оценки ожидаемого hit rate кэша размером $C$ при распределении Ципфа применяется аппроксимация Чи (Che's approximation) \cite{che_approx}:

\begin{equation}
h(n) \approx 1 - e^{-\lambda_n t_C(C)},
\label{eq:che}
\end{equation}

где $\lambda_n$ --- интенсивность запросов к $n$-й модели (определяется из~(\ref{eq:zipf}) умножением на общую интенсивность запросов), $t_C(C)$ --- характеристическое время кэша: при данном размере кэша $C$ объект с интенсивностью запросов $\lambda_n > 1/t_C$ будет храниться с высокой вероятностью.

Аппроксимация Чи позволяет оценить совокупный hit rate для заданного размера кэша без полного моделирования трасс запросов. Это полезно при планировании дискового ресурса нод: можно заранее оценить, какой размер кэша на ноде (например, 200~ГБ или 1~ТБ NVMe) даст достаточный hit rate для ожидаемого распределения запросов.

\subsection{Параллельная загрузка, выбор пира и prefetching}

\subsubsection*{Алгоритм выбора пира}

При наличии нескольких пиров, хранящих нужный чанк, система выбирает оптимальный узел для скачивания. Алгоритм выбора пира основан на следующих критериях, ранжированных по приоритету:

\begin{enumerate}[leftmargin=*]
  \item \textbf{Свежесть heartbeat}: пир с истёкшим или отсутствующим heartbeat немедленно исключается (\texttt{EXISTS node:<peer\_id>} = 0). Это защита от направления трафика к недоступным нодам.
  \item \textbf{Load score}: агент с меньшей вычислительной и сетевой нагрузкой предпочтительнее. Load score публикуется в heartbeat-записи Redis и обновляется каждые 10--30~с.
  \item \textbf{Локальность}: пир на той же стойке (rack) или в той же зоне доступности предпочтительнее за счёт меньшего RTT и более высокой пропускной способности внутрирячной сети.
  \item \textbf{Измеренный RTT}: периодические зондирующие запросы позволяют уточнить актуальную задержку до известных пиров.
\end{enumerate}

Оценочная функция пира: $\mathrm{score}(p) = \alpha \cdot \mathrm{latency\_ms}(p) + \beta \cdot \mathrm{load}(p)$, где $\alpha$ и $\beta$ --- весовые коэффициенты, настраиваемые в конфигурации агента. Выбирается пир с минимальным score.

\subsubsection*{Параллельная загрузка и закон Амдала}

Если несколько чанков модели отсутствуют в локальном кэше, и доступны несколько пиров, параллельная загрузка с разных пиров существенно сокращает общее время. Время параллельной загрузки \cite{zipf_law}:

\begin{equation}
T_{\mathrm{parallel}} = \frac{S}{B_{\mathrm{total}}} + \mathrm{RTT},
\label{eq:p2p_time}
\end{equation}

где $S$ --- суммарный объём загружаемых данных, $B_{\mathrm{total}} = \sum_{i=1}^{k} B_i$ --- суммарная полоса пропускания $k$ параллельных соединений, RTT --- задержка установления соединений (пренебрежимо мала в датацентре).

Ускорение параллельной загрузки ограничено последовательными частями процесса --- законом Амдала \cite{amdahl_law}:

\begin{equation}
S(n) = \frac{1}{(1-p) + p/n},
\label{eq:amdahl}
\end{equation}

где $p$ --- доля работы, поддающейся параллелизации (загрузка чанков), $(1-p)$ --- последовательная часть (координация через Redis, сборка результата, инициализация FUSE). При $p = 0.9$ и $n = 4$ пирах:
\[
S(4) = \frac{1}{0.1 + 0.225} = \frac{1}{0.325} \approx 3.1.
\]
При $n = 16$ пирах: $S(16) \approx 5.8$. Предел при $n \to \infty$: $S_{\max} = 1/(1-p) = 10$. Таким образом, увеличение числа параллельных источников сверх 8--10 даёт убывающую отдачу, и оптимальное число пиров для параллельной загрузки --- в диапазоне 4--8 для типичного датацентра с 10G-сетью.

\subsubsection*{Оптимальный размер чанка}

Выбор размера чанка представляет собой задачу оптимизации с двумя конкурирующими составляющими. Функция суммарных накладных расходов на единицу данных:

\begin{equation}
f(C) = \frac{T_{\mathrm{overhead}}}{C} + \frac{C}{B_{\mathrm{net}}},
\label{eq:chunk_cost}
\end{equation}

где $T_{\mathrm{overhead}}$ --- фиксированные накладные расходы на один чанк (переключение контекста FUSE, установка TCP-соединения, операции Redis), $B_{\mathrm{net}}$ --- пропускная способность сети. Минимум функции~(\ref{eq:chunk_cost}) достигается при значении, найденном из условия нулевой производной $df/dC = 0$:

\begin{equation}
C_{\mathrm{opt}} = \sqrt{T_{\mathrm{overhead}} \cdot B_{\mathrm{net}}}.
\label{eq:chunk_opt}
\end{equation}

При $T_{\mathrm{overhead}} = 1$~мс и $B_{\mathrm{net}} = 10$~Гбит/с = $1.25 \times 10^9$~Б/с:
\[
C_{\mathrm{opt}} = \sqrt{10^{-3} \times 1.25 \times 10^9} = \sqrt{1.25 \times 10^6} \approx 1118\,\text{КБ} \approx 1\,\text{МБ}.
\]
Реальные накладные расходы $T_{\mathrm{overhead}}$ включают несколько компонент и в зависимости от конфигурации могут составлять от единиц до десятков миллисекунд; при $T_{\mathrm{overhead}} = 10$~мс оптимальный размер чанка возрастает до $\approx 3.5$~МБ. Диапазон 8--64~МБ, применяемый на практике, соответствует консервативной оценке накладных расходов с запасом для обеспечения приемлемого time-to-first-byte.

\subsubsection*{Алгоритм prefetching}

При последовательном паттерне чтения (характерном для загрузки весов модели слой за слоем в vLLM) \texttt{storage-mounter} применяет алгоритм упреждающей загрузки. После обслуживания запроса чанка $k$ асинхронно запрашиваются чанки $k{+}1$, $k{+}2$, \ldots, $k{+}N$, где $N$ --- размер окна prefetching. К моменту, когда ML-runtime запрашивает чанк $k{+}1$, он уже находится в локальном буфере \texttt{storage-mounter}. Сетевая задержка скрыта за временем обработки предыдущего чанка GPU-ядром.

Алгоритм обнаруживает последовательный паттерн за $O(1)$: достаточно сравнить текущее смещение с ожидаемым (предыдущее смещение + размер чанка). При обнаружении паттерна запускаются $N$ параллельных горутин, каждая загружает один следующий чанк через UDS к \texttt{storage-agent}. Пространственная сложность буфера prefetching: $O(N \times C)$. Значение $N = 10$ при $C = 16$~МБ даёт буфер 160~МБ --- приемлемо для типичной ноды.

\subsection{Надёжность и высокая доступность}

\subsubsection*{Формула доступности при репликации}

Чанк считается доступным, если хотя бы одна нода, хранящая его, работоспособна. При факторе репликации $r$ и вероятности доступности каждой ноды $a$:

\begin{equation}
A = 1 - (1 - a)^r.
\label{eq:availability}
\end{equation}

При $a = 0.99$ и $r = 3$ репликах:
\[
A = 1 - (0.01)^3 = 1 - 10^{-6} = 0.999999.
\]

Вероятность одновременного отказа всех трёх реплик составляет $10^{-6}$ --- один случай на миллион «нодо-дней» \cite{ford2010availability}. Для сравнения: NFS-сервер с доступностью $a = 0.999$ (три девятки) при $r = 1$ (нет репликации) обеспечивает ту же $A = 0.999$. Добавление репликации в P2P-архитектуре даёт шесть девяток без изменения доступности отдельной ноды.

Репликация в предложенной архитектуре возникает органически: как только один агент загружает чанк из внешнего источника и регистрирует его в Redis, другие агенты могут найти его через \texttt{SMEMBERS} и при необходимости скачать копию. Администратор может задать минимальный фактор репликации для «горячих» моделей: агент при регистрации чанка проверяет, есть ли уже $r$ пиров в \texttt{:peers}, и если нет --- реплицирует самостоятельно.

\begin{figure}[htbp]
\centering
\includegraphics[width=0.78\textwidth]{chart_availability.png}
\caption{Простой в год в зависимости от коэффициента репликации чанков ($a = 0.99$). При $r = 3$ простой сокращается с 87.6~ч/год до 31.5~с/год~--- переход от двух к~четырём девяткам доступности.}
\label{fig:availability}
\end{figure}

\subsubsection*{TTL как защита от stale metadata}

TTL на Redis-записях автоматически устраняет stale metadata: если агент прекратил обновлять heartbeat, его записи исчезнут через TTL. Следующий запрос не будет направлен к этому агенту. Агент, обнаружив ошибку соединения с потенциальным пиром (\texttt{ECONNREFUSED}, \texttt{ETIMEDOUT}), немедленно исключает его из текущей операции и опционально выполняет \texttt{SREM chunk:<id>:peers <dead\_node>} для ускорения очистки реестра.

\subsubsection*{Fallback и degraded mode}

External fallback --- загрузка из HuggingFace или другого внешнего репозитория --- является конечной гарантией доступности. Пока внешний репозиторий доступен, система способна получить любой чанк, даже если кластерный кэш полностью пуст. Fallback не является оптимальным по времени (десятки секунд для больших чанков), но обеспечивает корректность операции.

При временной недоступности Redis агент переключается в деградированный режим: обслуживает запросы только из локального кэша или через fallback, без P2P-координации. Это предотвращает полный сбой системы при проблемах с Redis и соответствует принципу graceful degradation.

\subsubsection*{Eventual consistency для метаданных}

Redis-реестр обеспечивает eventual consistency для метаданных чанков. При добавлении новой ноды с чанками её записи становятся видимы другим агентам немедленно после успешного \texttt{SADD}. При удалении ноды её heartbeat-запись исчезнет через TTL. В промежутке между отказом ноды и истечением TTL её идентификатор присутствует в реестре, но попытка соединения вернёт ошибку, и агент корректно обработает её, выбрав следующий пир. Таким образом, eventual consistency в части метаданных не нарушает корректность алгоритма чтения --- лишь добавляет задержку в крайне редких пограничных случаях.

\subsection{Выводы по разделу}

Разработанная архитектура P2P/FUSE-системы кэширования устраняет все ключевые ограничения NFS-этапа и формирует многоуровневую систему доступа к данным ML-моделей. Данные распределены по локальным дискам worker-нод, что исключает единую точку отказа и узкое место по пропускной способности: каждая нода читает чанки из локального кэша или получает их от нескольких соседних нод параллельно.

Redis heartbeat обеспечивает $O(1)$-lookup активных агентов и доступных чанков --- принципиальное преимущество перед eventual-consistent Gossip (задержка $O(\log N)$ раундов) при cold start. FUSE-клиент предоставляет ML-runtime стандартный POSIX-интерфейс без изменения кода инференс-фреймворков. Lazy loading позволяет начать обработку модели до завершения загрузки всех весов, перекрывая фазы вычисления и ввода-вывода.

Ключевые аналитические результаты раздела:
\begin{itemize}[leftmargin=*]
  \item Трёхуровневая модель AMAT~(\ref{eq:amat}) при $H_l \approx 0.85$, $H_p \approx 0.90$ показывает, что более 98\% запросов обслуживаются без обращения к внешнему репозиторию.
  \item Закон Амдала~(\ref{eq:amdahl}) при $p = 0.9$ ограничивает практический предел ускорения параллельной загрузки: оптимальное число пиров --- 4--8, дальнейший рост числа параллельных источников даёт убывающую отдачу.
  \item Формула доступности~(\ref{eq:availability}) при $r = 3$ репликах на чанк и $a = 0.99$ даёт $A = 0.999999$ --- шесть девяток надёжности.
  \item Разделение UDS (локальная коммуникация) и TCP (межнодовая) оптимизирует как задержку (мкс-диапазон для UDS), так и надёжность (flow control TCP).
\end{itemize}

В совокупности перечисленные свойства формируют архитектуру, обеспечивающую существенное сокращение задержки cold start при повторных запусках одних и тех же ML-моделей в пределах кластера --- характерном сценарии для производственных систем инференса.
